% ═══════════════════════════════════════════════════════════════════════════
% LEHRERHANDREICHUNG DS-10: Beobachtungsstunde
% Kompakter Überblick mit allen 3 Varianten
% ═══════════════════════════════════════════════════════════════════════════

\documentclass[11pt, a4paper]{article}

\usepackage[utf8]{inputenc}
\usepackage[T1]{fontenc}
\usepackage[ngerman]{babel}
\usepackage{helvet}
\renewcommand{\familydefault}{\sfdefault}

\usepackage[margin=1.2cm, top=1cm, bottom=1cm]{geometry}
\usepackage[table]{xcolor}
\usepackage{array}
\usepackage{tabularx}
\usepackage{booktabs}
\usepackage{tikz}
\usepackage{tcolorbox}
\usepackage{enumitem}
\usepackage{amssymb}
\tcbuselibrary{skins}

% Farben
\definecolor{spanisch}{HTML}{c41e3a}
\definecolor{gruppeA}{HTML}{2a7a4b}
\definecolor{gruppeB}{HTML}{1565c0}
\definecolor{hellgrau}{HTML}{f0f0f0}
\definecolor{phase}{HTML}{666666}
\definecolor{wichtig}{HTML}{b35c00}
\definecolor{varianteA}{HTML}{c41e3a}
\definecolor{varianteB}{HTML}{1565c0}
\definecolor{varianteC}{HTML}{7b1fa2}

\pagestyle{empty}
\setlist{nosep, leftmargin=1.5em}

\begin{document}

% ═══════════════════════════════════════════════════════════════════════════
% HEADER
% ═══════════════════════════════════════════════════════════════════════════

\begin{tikzpicture}[remember picture, overlay]
\fill[spanisch] ([yshift=-1cm]current page.north west) rectangle ([yshift=-2.2cm]current page.north east);
\node[white, font=\Large\bfseries, anchor=west] at ([yshift=-1.6cm, xshift=1.2cm]current page.north west) {DS 10 -- BEOBACHTUNGSSTUNDE};
\node[white, font=\normalsize, anchor=east] at ([yshift=-1.6cm, xshift=-1.2cm]current page.north east) {90 Min | Di 03.02.2026};
\end{tikzpicture}

\vspace{1.2cm}

% ═══════════════════════════════════════════════════════════════════════════
% BEOBACHTUNGS-BOX
% ═══════════════════════════════════════════════════════════════════════════

\begin{tcolorbox}[
    colback=wichtig!15,
    colframe=wichtig,
    boxrule=1.5pt,
    top=4pt, bottom=4pt, left=6pt, right=6pt
]
\textbf{BEOBACHTUNG:} Nur 2. Hälfte (45 Min) wird beobachtet! 1. Hälfte = Vorbereitung. Zeige: Differenzierung, Sprechaktivität, Scaffolding.
\end{tcolorbox}

\vspace{0.2cm}

% ═══════════════════════════════════════════════════════════════════════════
% 3 VARIANTEN
% ═══════════════════════════════════════════════════════════════════════════

\begin{center}
{\large\bfseries Wähle EINE Variante:}
\end{center}

\vspace{0.2cm}

\begin{minipage}[t]{0.32\textwidth}
\begin{tcolorbox}[
    colback=varianteA!10,
    colframe=varianteA,
    title={\footnotesize\bfseries A: Tandem-Interview},
    fonttitle=\bfseries\color{white},
    colbacktitle=varianteA,
    boxrule=0.8pt,
    top=2pt, bottom=2pt, left=4pt, right=4pt
]
\footnotesize
\textbf{Konzept:} Homogene Paare → Gemischte Tandems → Präsentation

\textbf{Vorteile:}
\begin{itemize}
\item Explizite Differenzierung
\item Peer-Tutoring (B hilft A)
\item Transfer 1.→3. Person
\end{itemize}

\textbf{Risiko:} Mittel (Timing)

\textbf{Score:} 9.4/10
\end{tcolorbox}
\end{minipage}
\hfill
\begin{minipage}[t]{0.32\textwidth}
\begin{tcolorbox}[
    colback=varianteB!10,
    colframe=varianteB,
    title={\footnotesize\bfseries B: Speed-Dating},
    fonttitle=\bfseries\color{white},
    colbacktitle=varianteB,
    boxrule=0.8pt,
    top=2pt, bottom=2pt, left=4pt, right=4pt
]
\footnotesize
\textbf{Konzept:} 6 Runden à 3 Min, Gong, Wechsel

\textbf{Vorteile:}
\begin{itemize}
\item Timer steuert alles
\item Alle sprechen viel
\item Kein Schreiben während
\end{itemize}

\textbf{Risiko:} Niedrig

\textbf{Score:} 8.7/10
\end{tcolorbox}
\end{minipage}
\hfill
\begin{minipage}[t]{0.32\textwidth}
\begin{tcolorbox}[
    colback=varianteC!10,
    colframe=varianteC,
    title={\footnotesize\bfseries C: Galerierundgang},
    fonttitle=\bfseries\color{white},
    colbacktitle=varianteC,
    boxrule=0.8pt,
    top=2pt, bottom=2pt, left=4pt, right=4pt
]
\footnotesize
\textbf{Konzept:} Poster erstellen → Gallery Walk → Reflexion

\textbf{Vorteile:}
\begin{itemize}
\item Kreatives Produkt
\item Schreiben VOR Sprechen
\item Alle in Bewegung
\end{itemize}

\textbf{Risiko:} Niedrig-Mittel

\textbf{Score:} 9.0/10
\end{tcolorbox}
\end{minipage}

\vspace{0.3cm}
\hrule
\vspace{0.2cm}

% ═══════════════════════════════════════════════════════════════════════════
% VARIANTE A: VERLAUFSPLAN (EMPFOHLEN)
% ═══════════════════════════════════════════════════════════════════════════

\begin{tcolorbox}[
    colback=varianteA!5,
    colframe=varianteA,
    title={\bfseries VARIANTE A: Tandem-Interview (EMPFOHLEN)},
    fonttitle=\bfseries,
    boxrule=0.8pt,
    top=2pt, bottom=2pt, left=4pt, right=4pt
]
\footnotesize
\textbf{BEOBACHTETE 2. HÄLFTE (45 Min):}

\vspace{0.2cm}

\begin{tabularx}{\textwidth}{|c|c|X|c|}
\hline
\rowcolor{varianteA!20}
\textbf{Zeit} & \textbf{Phase} & \textbf{Aktivität} & \textbf{SF} \\
\hline
0--3 & Orientierung & Begrüßung, Ablauf erklären, Regeln: \textit{Solo español} & PL \\
\hline
\multicolumn{4}{|c|}{\cellcolor{hellgrau}\textbf{PHASE 1: Homogene Interviews (12 Min)}} \\
\hline
3--10 & Homogene Paare & A↔A, B↔B. Fragekarten verteilen. Fokus auf A-Paare. & PA \\
\hline
10--12 & Partnerwechsel & Kurzes 2. Interview mit neuem Partner gleicher Gruppe & PA \\
\hline
12--15 & Kurzes Feedback & 2--3 Antworten sammeln: \textit{¿Qué habéis aprendido?} & PL \\
\hline
\multicolumn{4}{|c|}{\cellcolor{hellgrau}\textbf{PHASE 2: Gemischte Tandems (18 Min) -- KERNPHASE}} \\
\hline
15--17 & Tandembildung & Tandems aufrufen, Ankreuz-Vorlage verteilen & -- \\
\hline
17--27 & B interviewt A & Timer 10 Min. B stellt Fragen, A antwortet. \textbf{B kreuzt an!} & PA \\
\hline
27--33 & A fragt B & Timer 6 Min. A nutzt einfache Fragen. & PA \\
\hline
\multicolumn{4}{|c|}{\cellcolor{hellgrau}\textbf{PHASE 3: Präsentation (9 Min) -- NUR 3 TANDEMS}} \\
\hline
33--42 & Vorstellung & E.H.+H.P., L.N.+A.A., V.S.+A.S. Je 3 Min. B stellt A vor. & PL \\
\hline
\multicolumn{4}{|c|}{\cellcolor{hellgrau}\textbf{PHASE 4: Abschluss (3 Min)}} \\
\hline
42--45 & Cierre & \textit{¡Muy bien! Ahora sabemos más sobre nuestros compañeros.} & PL \\
\hline
\end{tabularx}
\end{tcolorbox}

\vspace{0.2cm}

% ═══════════════════════════════════════════════════════════════════════════
% MATERIAL & TANDEMS
% ═══════════════════════════════════════════════════════════════════════════

\begin{minipage}[t]{0.48\textwidth}
\begin{tcolorbox}[
    colback=hellgrau,
    colframe=phase,
    title={\footnotesize\bfseries Material-Checkliste},
    fonttitle=\bfseries,
    boxrule=0.5pt,
    top=2pt, bottom=2pt, left=4pt, right=4pt
]
\footnotesize
\begin{itemize}
\item[$\square$] \textbf{MATERIAL-10-fragekarten} (A + B)
\item[$\square$] \textbf{MATERIAL-10-ankreuz-vorlage} (10x)
\item[$\square$] Timer (sichtbar für alle)
\item[$\square$] Tafel: Strukturhilfe 1.→3. Person
\end{itemize}

\vspace{0.1cm}
\textit{Optional für Variante C: A4-Blätter, Stifte, Klebestreifen}
\end{tcolorbox}
\end{minipage}
\hfill
\begin{minipage}[t]{0.48\textwidth}
\begin{tcolorbox}[
    colback=gruppeB!10,
    colframe=gruppeB,
    title={\footnotesize\bfseries Tandem-Zuordnung (Variante A)},
    fonttitle=\bfseries,
    boxrule=0.5pt,
    top=2pt, bottom=2pt, left=4pt, right=4pt
]
\footnotesize
\begin{tabular}{@{}lll@{}}
\textbf{B (interviewt)} & \textbf{A (antwortet)} & \textbf{Hinweis} \\
\hline
E.H. (TOP) & H.P. & E.H. als Vorbild \\
L.N. (proaktiv) & A.A. & L.N. ermutigt \\
V.S. (konstant) & A.S. & ruhiges Tandem \\
C.P. (ruhig) & L.K. & beide konzentriert \\
{[S9]} & {[S10]} & je nach Anwesenheit \\
\end{tabular}
\end{tcolorbox}
\end{minipage}

\vspace{0.2cm}

% ═══════════════════════════════════════════════════════════════════════════
% STRUKTURHILFE
% ═══════════════════════════════════════════════════════════════════════════

\begin{tcolorbox}[
    colback=white,
    colframe=spanisch,
    title={\footnotesize\bfseries Strukturhilfe (Tafelanschrieb)},
    fonttitle=\bfseries,
    boxrule=0.5pt,
    top=2pt, bottom=2pt, left=4pt, right=4pt
]
\footnotesize
\textbf{Umformung 1. Person → 3. Person:}
\hspace{1cm}
\textit{Me llamo María} → \textbf{Se llama} María
\hspace{0.5cm}|\hspace{0.5cm}
\textit{Soy simpática} → \textbf{Es} simpática
\hspace{0.5cm}|\hspace{0.5cm}
\textit{Me gusta bailar} → \textbf{Le gusta} bailar
\end{tcolorbox}

\vspace{0.2cm}

% ═══════════════════════════════════════════════════════════════════════════
% SCHWIERIGKEITEN
% ═══════════════════════════════════════════════════════════════════════════

\begin{tcolorbox}[
    colback=white,
    colframe=wichtig,
    title={\footnotesize\bfseries Erwartete Schwierigkeiten},
    fonttitle=\bfseries,
    boxrule=0.5pt,
    top=2pt, bottom=2pt, left=4pt, right=4pt
]
\footnotesize
\begin{tabularx}{\textwidth}{@{}p{4cm}X@{}}
\textbf{A versteht Fragen von B nicht} & B vereinfacht, L gibt Beispiel, Fragekarte A als Backup \\
\textbf{Umformung 1.→3. Person schwierig} & Strukturhilfe an Tafel zeigen, B liest von Vorlage ab \\
\textbf{Nervosität wegen Beobachtung} & Routine durch 1. Hälfte, bekannte Strukturen, positives Feedback \\
\textbf{Zeitproblem Phase 2} & Rollenwechsel kürzen (3 statt 6 Min) oder nur B→A \\
\end{tabularx}
\end{tcolorbox}

\end{document}
